\setcounter{section}{8}






  \zwischenueberschrift{Übungsaufgaben}

\inputaufgabe
{}
{





Bestimme die möglichen Durchschnitte von zwei zueinander senkrechten Zylindern.





}
{} {}

\inputaufgabe
{}
{





Erstelle eine möglichst einfache Gleichung für das mechanische System, das durch die $X$-Achse, die
\zusatzklammer {verschobene} {} {}
Parabel

\mavergleichskette
{\vergleichskette
{Y
}
{ = }{ X^2+1
}
{  }{
}
{    }{
}
{   }{
}
}
{}{}{}
und den Abstand $2$ gegeben ist.





}
{} {}

\inputaufgabe
{}
{





Es seien

\mavergleichskette
{\vergleichskette
{ H_1,H_2
}
{ \in }{ K[X]
}
{  }{
}
{    }{
}
{   }{
}
}
{}{}{}
Polynome in einer Variablen und

\mavergleichskette
{\vergleichskette
{ C_1
}
{ = }{ V(Y-H_1)
}
{  }{
}
{    }{
}
{   }{
}
}
{}{}{}
und

\mavergleichskette
{\vergleichskette
{ C_2
}
{ = }{ V(Y-H_2)
}
{  }{
}
{    }{
}
{   }{
}
}
{}{}{}
die zugehörigen
\definitionsverweis {Graphen}{}{}
im
\mathl{{\mathbb A}^{2}_{ K }}{.} Zeige, dass man das zugehörige
\definitionsverweis {mechanische System}{}{}
\zusatzklammer {zu einem bestimmten Abstand $d$} {} {}
mit zwei Variablen beschreiben kann.





}
{} {}

\inputaufgabe
{}
{





Bestimme Gleichungen für das
\definitionsverweis {mechanische System}{}{,}
das durch den Einheitskreis, die $x$-Achse und den Abstand $1$ gegeben ist. Was sind die
\definitionsverweis {irreduziblen Komponenten}{}{}
des Systems?





}
{} {}

\inputaufgabe
{}
{





Es sei

\mavergleichskette
{\vergleichskette
{M
}
{ \subset }{ { {\mathbb A}_{ \R }^{  4   } }
}
{  }{
}
{    }{
}
{   }{
}
}
{}{}{}
ein durch die beiden Bahnen
\mathkor {} {B_1} {und} {B_2} {}
und den Abstand $d$ gegebenes
\definitionsverweis {mechanisches System}{}{.}
Zeige, dass es eine natürliche
\definitionsverweis {injektive}{}{}
Abbildung
\maabbdisp {} { M } { B_1 \times B_2
} {}
gibt.





}
{} {}

\inputaufgabe
{}
{





Es sei

\mavergleichskettedisp
{\vergleichskette
{B
}
{ =} { V(G)
}
{ \subset} { {\mathbb A}^{2}_{\R}
}
{ } {
}
{ } {
}
}
{}{}{}
eine Bahn, die wir als
\definitionsverweis {mechanisches System}{}{}
$M$ in dem Sinne auffassen, dass die beiden Punkte mit dem Abstand $d$ auf dieser einen Bahn liegen müssen. Zeige, dass es eine natürliche
\definitionsverweis {fixpunktfreie}{}{}
\definitionsverweis {bijektive}{}{}
Abbildung
\maabbdisp {} { M } { M
} {}
gibt.





}
{} {}

\inputaufgabe
{}
{





Bestimme Gleichungen für das
\definitionsverweis {mechanische System}{}{,}
das durch die $x$-Achse
\zusatzklammer {als gemeinsame Bahn} {} {}
und den Abstand $1$ gegeben ist.





}
{} {}

\inputaufgabe
{}
{





Bestimme Gleichungen für das
\definitionsverweis {mechanische System}{}{,}
das durch den
\definitionsverweis {Einheitskreis}{}{}
\zusatzklammer {als gemeinsame Bahn} {} {}
und den Abstand $2$ gegeben ist.





}
{} {}

\inputaufgabegibtloesung
{}
{





Wir betrachten das
\definitionsverweis {mechanische System}{}{,}
das durch den Einheitskreis und die dazu tangentiale Gerade durch
\mathl{(0,1)}{} mit dem Koppelungsabstand

\mavergleichskette
{\vergleichskette
{ d
}
{ = }{ 2
}
{  }{
}
{    }{
}
{   }{
}
}
{}{}{}
definiert ist. Zeige, dass man dieses System mit zwei Variablen beschreiben kann.





}
{} {}

\inputaufgabe
{}
{





Wir betrachten das
\definitionsverweis {mechanische System}{}{,}
das durch den Einheitskreis und die dazu tangentiale Gerade durch
\mathl{(0,1)}{} mit dem Koppelungsabstand

\mavergleichskette
{\vergleichskette
{ d
}
{ = }{ 2
}
{  }{
}
{    }{
}
{   }{
}
}
{}{}{}
definiert ist, siehe

Beispiel 8.5
.
Zeige, dass dieses System
\definitionsverweis {irreduzibel}{}{}
ist. Wende dazu

Aufgabe 7.24

auf die Ringerweiterung

\mavergleichskettedisp
{\vergleichskette
{ R
}
{ =} { \R[X_2,Y_2]/ \left(  X_2^2+Y_2^2-1  \right)
}
{ \subseteq} { R [X_1]/ \left(  X_1^2- 2X_1X_2-2Y_2 -2    \right)
}
{ } {
}
{ } {
}
}
{}{}{}
an.





}
{} {}

\inputaufgabe
{}
{





Bestimme Gleichungen für das
\definitionsverweis {mechanische System}{}{,}
das durch die $x$-Achse, die $y$-Achse und den Abstand $1$ gegeben ist.





}
{} {}

\inputaufgabe
{}
{





Bestimme Gleichungen für das
\definitionsverweis {mechanische System}{}{,}
das durch das Achsenkreuz
\zusatzklammer {als gemeinsame Bahn} {} {}
und den Abstand $1$ gegeben ist.





}
{} {}

\inputaufgabe
{}
{





Es seien

\mavergleichskettedisp
{\vergleichskette
{B_1 = V(F_1), B_2 = V(F_2)
}
{ \subseteq} { {\mathbb A}^{2}_{\R}
}
{ } {
}
{ } {
}
{ } {
}
}
{}{}{}
zwei Bahnen, und es sei ein Abstand $d$ fixiert. Vergleiche das
\definitionsverweis {mechanische System}{}{}
$M$ zu diesen Bahnen mit dem System $N$, das zu der einen Bahn
\mathl{B_1 \cup B_2}{} gehört. Zeige, dass es zwei natürliche injektive Abbildungen
\maabbdisp {} { M } { N
} {}
gibt. Es sei $L_i$ das mechanische System das zu $B_i$ als alleiniger Bahn gehört. Zeige, dass es eine natürliche surjektive Abbildung der Form
\maabbdisp {} {L_1 \uplus L_2 \uplus M \uplus M} {N
} {}
gibt.





}
{} {}

\inputaufgabe
{}
{





Es sei

\mavergleichskette
{\vergleichskette
{ M
}
{ \subset }{ { {\mathbb A}_{ \R }^{  4   } }
}
{  }{
}
{    }{
}
{   }{}
}
{}{}{}
ein
\definitionsverweis {mechanisches System}{}{.}
Zeige, dass durch
\maabbeledisp {} { M } { S^1(d)
} {(P_1,P_2)} { P_1-P_2
} {,}
eine Abbildung des Systems auf den Kreis mit Mittelpunkt
\mathl{(0,0)}{} und Radius $d$ gegeben ist. Was bedeutet die Surjektivität dieser Abbildung? Kann diese Abbildung nur endlich viele Bildpunkte besitzen?





}
{} {}

\inputaufgabe
{}
{





Wir betrachten das
\definitionsverweis {mechanische System}{}{}

\mavergleichskette
{\vergleichskette
{M
}
{ \subset }{ { {\mathbb A}_{ \R }^{  4   } }
}
{  }{
}
{    }{
}
{   }{
}
}
{}{}{}
zu zwei sich kreuzenden Geraden. Zeige, dass die Abbildung aus

Aufgabe 8.14

bijektiv ist.





}
{} {}

\inputaufgabe
{}
{





Wir betrachten das
\definitionsverweis {mechanische System}{}{}

\mavergleichskette
{\vergleichskette
{M
}
{ \subset }{ { {\mathbb A}_{ \R }^{  4   } }
}
{  }{
}
{    }{
}
{   }{
}
}
{}{}{}
zum Einheitskreis, zum Kreis mit Mittelpunkt
\mathl{(0,2)}{} und Radius $4$ und zum Koppelungsabstand $5$. Zeige, dass die Abbildung aus

Aufgabe 8.14

nicht surjektiv ist. Was ist das Bild?





}
{} {}

In den folgenden Aufgaben betrachten wir die Nullstellenmenge

\mavergleichskettedisp
{\vergleichskette
{V
}
{ =} { V(f_1  , \ldots , f_k)
}
{ \subset} {  { {\mathbb A}_{ \R }^{ n  } }
}
{ } {
}
{ } {
}
}
{}{}{}
differentialgeometrisch. Wir erinnern an die folgende Definition von einem regulären Punkt zu einer differenzierbaren Abbildung aus der Analysis 2 Vorlesung. Dieser Begriff ist im jetzigen Kontext auf die Abbildung
\maabbdisp {\varphi  = (f_1  , \ldots , f_k)} {{ {\mathbb A}_{ \R }^{ n  } }  } {{ {\mathbb A}_{ \R }^{ k  } }
} {}
in einem Punkt $P \in V$ anzuwenden.







Es seien
\mathkor {} {V} {und} {W} {}
\definitionsverweis {endlichdimensionale}{}{}
\definitionsverweis {reelle Vektorräume}{}{,}
sei

\mavergleichskette
{\vergleichskette
{G
}
{ \subseteq }{ V
}
{  }{
}
{    }{
}
{   }{
}
}
{}{}{}
\definitionsverweis {offen}{}{,}
sei

\mavergleichskette
{\vergleichskette
{P
}
{ \in }{ G
}
{  }{
}
{    }{
}
{   }{
}
}
{}{}{}
und sei
\maabbdisp {\varphi} { G } { W
} {}
eine in $P$
\definitionsverweis {differenzierbare Abbildung}{}{.}
Dann heißt $P$ ein \definitionswort {regulärer Punkt}{} von $\varphi$, wenn

\mavergleichskettedisp
{\vergleichskette
{ \operatorname{rang} \,  \left(D\varphi\right)_{P}
}
{ =} {  {\min   \left(   \dim_{  } \left(  V  \right)  ,  \dim_{  } \left(  W  \right)   \right)   }
}
{ } {
}
{ } {
}
{ } {
}
}
{}{}{}
ist. Andernfalls heißt $P$ ein \definitionswort {kritischer Punkt}{} oder ein \definitionswort {singulärer Punkt}{.}







\inputaufgabegibtloesung
{}
{





Wir betrachten das
\definitionsverweis {mechanische System}{}{,}
das durch die $x$-Achse und den Kreis mit Radius $1$ und Mittelpunkt
\mathl{(0,2)}{} gegeben ist. Der Koppelungsabstand sei

\mavergleichskette
{\vergleichskette
{ d
}
{ > }{ 0
}
{  }{
}
{    }{
}
{   }{
}
}
{}{}{.}
\aufzaehlungdreiabc{Erstelle die Gleichungen, die dieses System beschreiben.
}{Bestimme, für welche $d$ das System in jedem Punkt
\definitionsverweis {regulär}{}{}
ist.
}{Bestimme die kritischen Punkte in Abhängigkeit von $d$. Wie kann man diese Punkte als Eigenschaft des mechanischen Systems erklären?
}





}
{} {}

\inputaufgabe
{}
{





Wir betrachten das
\definitionsverweis {mechanische System}{}{,}
das durch die $x$-Achse und den Kreis mit Radius $1$ und Mittelpunkt
\mathl{(0,2)}{} gegeben ist. Der Koppelungsabstand sei $d>0$.
\aufzaehlungzwei {Begründe anschaulich und zeige, dass dieses mechanische System zu

\mavergleichskette
{\vergleichskette
{d
}
{ > }{ 3
}
{  }{
}
{    }{
}
{   }{
}
}
{}{}{}
nicht
\zusatzklammer {in der reellen Topologie} {} {}
\definitionsverweis {wegzusammenhängend}{}{}
ist.
} {Begründe anschaulich und zeige, dass dieses mechanische System zu

\mavergleichskette
{\vergleichskette
{ 1
}
{ \leq }{d
}
{ \leq }{ 3
}
{    }{
}
{   }{
}
}
{}{}{}
\definitionsverweis {wegzusammenhängend}{}{}
ist.
}





}
{} {}





  \zwischenueberschrift{Aufgaben zum Abgeben}

\inputaufgabe
{6 (2+2+2)}
{





Wir betrachten das
\definitionsverweis {mechanische System}{}{,}
das durch die $x$-Achse, die Parabel
\mathl{V(Y-X^2)}{} und den Koppelungsabstand $1$ gegeben ist.
\aufzaehlungdrei{Bestimme Gleichungen
\zusatzklammer {in möglichst wenigen Variablen} {} {}
für das mechanische System.
}{Besitzt das System kritische Punkte?
}{Bestimme die Gleichung für den Bewegungsvorgang zum Mittelpunkt der Verbindungsstange.
}





}
{} {}

\inputaufgabe
{2}
{





Bestimme Gleichungen für das
\definitionsverweis {mechanische System}{}{,}
das durch den
\definitionsverweis {Einheitskreis}{}{}
\zusatzklammer {als gemeinsame Bahn} {} {}
und den Abstand $1$ gegeben ist.





}
{} {}

\inputaufgabe
{6 (2+2+2)}
{





Wir betrachten das
\definitionsverweis {mechanische System}{}{,}
das durch die Parabel
\mathl{V(Y-X^2)}{,} den Kreis mit Mittelpunkt
\mathl{(0,2)}{} und Radius $1$ und den Koppelungsabstand $1$ gegeben ist.
\aufzaehlungdrei{Bestimme Gleichungen
\zusatzklammer {in möglichst wenigen Variablen} {} {}
für dieses mechanische System.
}{Bestimme die
\definitionsverweis {Zusammenhangskomponenten}{}{}
des Systems in der metrischen Topologie.
}{Bestimme die Zusammenhangskomponenten des Systems in der
\definitionsverweis {Zaris\-ki-Topologie}{}{.}
}





}
{} {}

\inputaufgabe
{6 (2+4)}
{





Wir betrachten das
\definitionsverweis {mechanische System}{}{,}
das durch die $x$-Achse und den Kreis mit Radius $1$ und Mittelpunkt
\mathl{(0,2)}{} gegeben ist. Der Koppelungsabstand sei

\mavergleichskette
{\vergleichskette
{ d
}
{ > }{ 0
}
{  }{
}
{    }{
}
{   }{
}
}
{}{}{.}
Wir knüpfen an

Aufgabe 8.17

an.
\aufzaehlungzweiabc{Eliminiere die Variable $y_1$ aus den Gleichungen des Systems.
}{Bestimme mit der einen beschreibenden Gleichung des Systems in den Variablen
\mathkor {} {x_1} {und} {x_2} {,}
für welche $d$ das System in jedem Punkt
\definitionsverweis {regulär}{}{}
ist.
}





}
{} {}

\inputaufgabe
{6}
{





Es sei

\mavergleichskette
{\vergleichskette
{ C
}
{ \subseteq }{ { {\mathbb A}_{ \R }^{  3   } }
}
{  }{
}
{    }{
}
{   }{
}
}
{}{}{}
der Durchschnitt von zwei Zylindern
\zusatzklammer {um die $x$- bzw. um die $y$-Achse} {} {}
mit Radius $1$
\zusatzklammer {$C$ ist also die Vereinigung von zwei Ellipsen} {} {.}
Wir betrachten die durch einen Vektor

\mavergleichskette
{\vergleichskette
{ v
}
{ = }{ (a,b,c)
}
{ \neq }{ 0
}
{    }{
}
{   }{
}
}
{}{}{}
definierte
\definitionsverweis {senkrechte Projektion}{}{}
\maabbdisp {p_v} { { {\mathbb A}_{ \R }^{  3   } } } { {\mathbb A}^{2}_{\R}
} {.}
Man charakterisiere, in Abhängigkeit von
\mathl{a,b,c}{,} die möglichen Bilder unter diesen Projektionen.





}
{} {}

\inputaufgabe
{3}
{





Betrachte die Abbildung
\maabbeledisp {} { {\mathbb A}^{2}_{ K } } { {\mathbb A}^{2}_{ K }
} { (x,y) } {  \left(  x^2   , \,   y^2                   \right) = (u,v)
} {.}
Wie sieht das Bild der Ebene und wie das Bild des Einheitskreises unter dieser Abbildung für

\mavergleichskette
{\vergleichskette
{ K
}
{ = }{ \R
}
{  }{
}
{    }{
}
{   }{
}
}
{}{}{}
und wie für

\mavergleichskette
{\vergleichskette
{ K
}
{ = }{ {\mathbb C}
}
{  }{
}
{    }{
}
{   }{
}
}
{}{}{}
aus? Im reellen Fall, wenn der Kreis einmal durchlaufen wird, wie oft wird das Bild durchlaufen?





}
{} {}
